% Pillar-7 (ZVF theory + adaptive-G controller) introduction + related work
\section{Introduction}
\label{sec:p7-intro}
GRPO~\citep{shao2024deepseekmath, guo2025deepseekr1} centers rewards within a
group of $G$ completions per prompt. Under binary, verifiable rewards, any
group whose completions are all correct or all incorrect has zero centered
reward contrast in the stated GRPO term; this does not imply that every
total-gradient contribution vanishes. The fraction of such degenerate groups---the Zero-Variance
Fraction (ZVF)---was the subject of our companion Pillar-2 paper, which
measured it across libraries, model scales, and training phases and
deliberately stopped at description: ZVF tracks collapse, drifts upward, and
is mechanically entangled with accuracy and group size, so Pillar~2 declined
to promote it to a predictive or causal statistic. A companion
theorem-level treatment of the same quantities --- identifiability, CI
calibration, the reliability budget, and the limits of the proposed
signal-per-rollout objective --- appears in the ZVF-program theory paper
\emph{``Calibrating the Zero-Variance Fraction in Group-Relative RL''}
(\texttt{zvf-program/theory/zvf\_theory.tex}); this paper is canonical for
the controller engineering, that one for the theory.

This paper takes the next two steps. First, \emph{prediction}: we show that a
minimal binomial model of group degeneracy---usable signal
$S = p(1-p)\,(1 - h_G(p))$ with $h_G(p) = p^G + (1-p)^G$---predicts, in
closed form, three regularities that Pillar~2 could only observe: the
expected ZVF indicator equals $p^G + (1-p)^G$ (T1), larger $G$ lowers ZVF
(T2), and gradient magnitude tracks the contrast term $p(1-p)$ (T3). Each
theorem is paired with an empirical check drawn from a 368-run
Weights~\&~Biases audit (which counts \emph{all} logged runs---training,
evaluation, and short accuracy-swept probes---and is thus a superset of the
benchmark's 70+ core \emph{training} runs), a model~$\times$~group-size grid, and an
open-backward-pass gradient probe. Second, \emph{intervention}: if the model
is right, then ZVF is not fate but a control target. We build a
\texttt{zvf-triage} callback that watches per-step ZVF and an adaptive-$G$
controller that escalates group size when the signal starves, and we audit it
head-to-head against fixed recipes (GRPO, Dr.~GRPO, DAPO-style dynamic
sampling) in a single open, auditable reimplementation.

The theory also explains why ZVF alone cannot drive a controller. Because
$h_G(p)$ is symmetric in $p \leftrightarrow 1-p$, ZVF assigns the same value
to a mastered prompt ($p \to 1$) and a hopeless one ($p \to 0$); at
population scale this produces the U-shape we measure across seven models.
Worse, ZVF is an \emph{existence} statistic: adding reward jitter
of order $10^{-4}$ drives measured ZVF to zero while shifting the measured
within-group contrast by less than $10^{-6}$. We therefore promote the
\emph{pairwise-contrast density}
$\mathrm{PCD} = \frac{G-1}{G}\,\E[p(1-p)]$---the expected within-group reward
variance---as the magnitude-aware statistic a controller should consume, and
verify empirically that PCD is invariant to the jitter that fools ZVF.

Our contributions are:
(i) a closed-form signal-starvation model of binary-reward GRPO with three
theorems, each paired with an empirical check;
(ii) the PCD statistic and a micro-jitter falsification separating it from
ZVF;
(iii) a toy-scale validation that gradient norms track $p(1-p)$
($\mathrm{corr} = +0.71$; directional only);
(iv) an operational \texttt{zvf-triage} callback and adaptive-$G$ controller,
audited in a four-arm open-trainer comparison against GRPO, Dr.~GRPO, and
DAPO-style dynamic sampling, with the accuracy-versus-compute trade stated
explicitly; and
(v) design rules for when and how to escalate $G$.
Throughout we keep the Pillar-2 discipline: toy-scale results are marked
directional, and every interventional claim is scoped to the single task and
small $n$ on which it was measured. The word \emph{controller} names the
implemented policy, not an established improvement: the decisive static-
$G{=}16$ matched-token bakeoff remains unrun.

\subsection{Related Work}
\label{sec:p7-related}
The degenerate-group problem is increasingly attacked at the sampler. DAPO
\citep{yu2025dapo} filters out zero-variance groups by resampling until the
batch carries signal (dynamic sampling); GRESO \citep{zhang2026greso} makes
the rollout decision itself selective to avoid paying for uninformative
prompts; RL-ZVP/AERO \citep{le2025rlzvp} instead extracts signal \emph{from}
zero-variance groups via entropy-guided advantage shaping. Our adaptive-$G$
controller is a third point in this design space: rather than discarding or
reshaping degenerate groups, it changes the group size that generates them,
guided by the closed-form $h_G(p)$. Dr.~GRPO \citep{tong2025drgrpo} and GSPO
\citep{qwen2025gspo} modify normalization and the importance-ratio level
respectively, and serve as fixed-recipe baselines in our audit. The
contrast term $p(1-p)$ at the heart of our model echoes the reduction of
group-relative RL to pairwise preference objectives
\citep{wu2025grpodpo, rafailov2024direct}: a group is exactly as informative
as the correct/incorrect pairs it contains. Empirical scaling studies of
GRPO post-training \citep{tan2025scalingrl} treat group size as a budget
knob; we treat it as a feedback-controlled variable. All tasks are built on
GSM8K-style verifiable rewards \citep{cobbe2021training}, and the benchmark
substrate, statistics, and seed policy are shared with the companion pillars
(Papers~1--4 of this series) and described in
Sections~\ref{sec:benchmark}--\ref{sec:setup}.
